% ==============================================================================
% Chapter 6: Problems and Comments
% ==============================================================================
\chapter{Problems and Comments}

\section{Student}
\subsection{Problems}
\begin{enumerate}[label=\arabic*.]
    \item I often felt sleepy in the morning and afternoon. I woke up at 6 a.m. and returned home around 8:30 to 9:00 p.m., so the long days reduced my energy and focus.
    \item My task management was disorganized. I often estimated task durations incorrectly and sometimes planned the order of tasks wrongly, so my lead advised me to discuss task plans with my mentor before creating tasks.
    \item I had used C\# for less than a year, so I often could not tell whether my implementation was good, professional, and scalable enough.
    \item I started game development only about one year ago, with roughly 7 months of actual development experience, so my game design skills are still developing. This showed in the pitching phase: my first idea was dropped because I did not consider the target audience, and my second idea needed to be redesigned because its difficulty curve was unusual.
    \item Progress was much slower than in my solo projects because every change went through strict pull request review. According to my company timesheet, implementing player movement took 7 hours 15 minutes in the prototype, but 2 days 2 hours 45 minutes for the controller movement logic and 6 hours for the movement state in production. Counting a working day as 8 hours, that is about 24 hours 45 minutes, roughly 3.4 times as long. I reviewed my own pull requests before submitting them, but with limited experience I could not always spot the issues that reviewers found.
    \item Tasks had to be done in order, so I could not move to another system while waiting for a pull request to merge. Instead, my branches were stacked on top of each other, and a mistake could leave a branch with the wrong parent.
    \item My Git knowledge had gaps. I had never used cherry-picking, and I sometimes needed to ask for help with the Git flow.
    \item I was sometimes hesitant to ask questions. I wanted to try solving a problem myself first, which sometimes worked and sometimes broke things, and I did not want to waste the time of the team as an intern.
    \item My performance review noted that I was slow to act on feedback. I understood each comment when I received it, but I did not always change my approach quickly enough, so the same kind of issue could appear again.
    \item My performance review also noted that I sometimes solved problems in ways that went outside the design the team had agreed on. When a problem did not fit the agreed design, I tried to fix it on my own instead of raising it first.
\item My communication was below what the team expected. I sometimes described problems unclearly or reported them too late, which meant my mentor had to help resolve them after they had already grown. I now try to raise problems early, even before I have a solution.
\end{enumerate}

\subsection{Comments}
\begin{enumerate}[label=\arabic*.]
    \item I should plan my sleep around the early start and late return, for example by going to bed earlier, so that I have enough energy during the day.
    \item I should discuss the task order and estimates with my mentor before creating tasks, break tasks into smaller pieces, add buffer time, and compare each estimate with the actual time spent so that my estimates improve over time.
    \item I should ask mentors about uncertain design decisions before implementing them, so that I do not have to rely only on my own judgment.
    \item I should study game design more deliberately, for example by analyzing how finished games handle difficulty, pacing, and their target audience, and by getting feedback on a pitch from the team and from playtesters before investing time in it.
    \newpage
    \item I should turn the review comments I receive into a checklist of recurring issues, and read the review comments on my teammates' pull requests to learn what reviewers look for. This builds the experience that I need to spot issues in my own pull requests.
    \item I should check the parent branch before creating each new branch, and plan the order of tasks in advance so that the branch stack stays short.
    \item I should study the Git commands I have not used, such as cherry-pick and rebase, and practice them in a test repository before using them on real work.
    \item I should set a time limit for trying a problem alone, for example 30 minutes, and ask after that. Asking early costs less of the team's time than fixing a large mistake later.
    \item When I receive feedback, I should write down what I will change, apply it to my next piece of work immediately, and check it against the earlier comment, so that I act on feedback quickly and do not repeat it.
    \item When a problem does not fit the agreed design, I should describe the problem and my proposed solution to my lead or mentor before implementing it, so that any change to the design is a team decision.
    \item I should give short progress updates on a regular schedule, and report blockers and changes in my plan as soon as I notice them, instead of waiting until I have a complete answer.
\end{enumerate}

\section{Company}
\subsection{Problems}
\begin{enumerate}[label=\arabic*.]
    \item Some of the coding conventions on the company wiki were out of date, so they did not always match what reviewers expected.
\end{enumerate}

\subsection{Comments}
\begin{enumerate}[label=\arabic*.]
    \item The wiki could be reviewed and updated periodically, so that new developers can rely on it as the single source for conventions.
\end{enumerate}

\section{University}
\subsection{Problems}
\begin{enumerate}[label=\arabic*.]
    \item The professor's workplace visit was scheduled late and with a very short deadline. As of 8 November 2026, I had not been contacted, while many of my classmates already had appointments or had already been visited. The university staff contacted me on 12 November at 10:00 a.m. to say that the professor was available on 24 November at 10:00, and asked me to confirm by 15:00 on the same day. This gave me only five hours to confirm a date that also depended on my company schedule, and the short deadline fell on me even though the timing of the contact was not in my control.
\end{enumerate}

\subsection{Comments}
\begin{enumerate}[label=\arabic*.]
    \item It would help if the department announced the expected timing of the visits to all students at the start of the internship, and contacted each student early enough to allow time to check with the company supervisor. A deadline of at least one or two working days for confirming a date would also be more reasonable than a few hours.
\end{enumerate}